\begin{lemma}
\label{lemma-structure-torsion-D-module-regular}
\begin{reference}
Special case of \cite[Theorem 2.4]{Lyubeznik}
\end{reference}
Let $k$ be a field of characteristic $0$. Let $d \geq 1$.
Let $A = k[[x_1, \ldots, x_d]]$ with maximal ideal $\mathfrak m$.
Let $M$ be an $\mathfrak m$-power torsion $A$-module endowed with
additive operators $D_1, \ldots, D_d$ satisfying the leibniz rule
$$
D_i(fz) = \partial_i(f) z + f D_i(z)
$$
for $f \in A$ and $z \in M$. Here $\partial_i$ is
differentiation with respect to $x_i$.
Then $M$ is isomorphic to a direct sum
of copies of the injective hull $E$ of $k$.
\end{lemma}

\begin{proof}
Choose a set $J$ and an isomorphism $M[\mathfrak m] \to \bigoplus_{j \in J} k$.
Since $\bigoplus_{j \in J} E$ is injective
(Dualizing Complexes, Lemma \ref{dualizing-lemma-sum-injective-modules})
we can extend this isomorphism to an $A$-module homomorphism
$\varphi : M \to \bigoplus_{j \in J} E$.
We claim that $\varphi$ is an isomorphism, i.e., bijective.

\medskip\noindent
Injective. Let $z \in M$ be nonzero. Since $M$ is $\mathfrak m$-power torsion
we can choose an element $f \in A$ such that $fz \in M[\mathfrak m]$ and
$fz \not = 0$. Then $\varphi(fz) = f\varphi(z)$ is nonzero, hence
$\varphi(z)$ is nonzero.

\medskip\noindent
Surjective. Let $z \in M$. Then $x_1^n z = 0$ for some $n \geq 0$.
We will prove that $z \in x_1M$ by induction on $n$.
If $n = 0$, then $z = 0$ and the result is true.
If $n > 0$, then applying $D_1$ we find $0 = n x_1^{n - 1} z + x_1^nD_1(z)$.
Hence $x_1^{n - 1}(nz + x_1D_1(z)) = 0$. By induction we get
$nz + x_1D_1(z) \in x_1M$. Since $n$ is invertible, we conclude
$z \in x_1M$. Thus we see that $M$ is $x_1$-divisible.
If $\varphi$ is not surjective, then we can choose
$e \in \bigoplus_{j \in J} E$ not in $M$.
Arguing as above we may assume $\mathfrak m e \subset M$,
in particular $x_1 e \in M$. There exists an element
$z_1 \in M$ with $x_1 z_1 = x_1 e$. Hence
$x_1(z_1 - e) = 0$. Replacing $e$ by $e - z_1$
we may assume $e$ is annihilated by $x_1$.
Thus it suffices to prove that
$$
\varphi[x_1] :
M[x_1]
\longrightarrow
\left(\bigoplus\nolimits_{j \in J} E\right)[x_1] =
\bigoplus\nolimits_{j \in J} E[x_1]
$$
is surjective. If $d = 1$, this is true by construction of $\varphi$.
If $d > 1$, then we observe that $E[x_1]$ is the injective hull
of the residue field of $k[[x_2, \ldots, x_d]]$, see
Dualizing Complexes, Lemma \ref{dualizing-lemma-quotient}.
Observe that $M[x_1]$ as a module over $k[[x_2, \ldots, x_d]]$
is $\mathfrak m/(x_1)$-power torsion and comes
equipped with operators $D_2, \ldots, D_d$ satisfying
the displayed Leibniz rule.
Thus by induction on $d$ we conclude that $\varphi[x_1]$
is surjective as desired.
\end{proof}